\chapter{绪论}
\section{引言}
研究背景

本文主要研究水下目标在复杂物理环境（如极化带效应、涡流效应等）作用下所产生的电荷的复杂时空分布特性，通过构建一套高敏感度与选择性的电荷传感器阵列系统，实现对目标的实时远距离探测与分析，最终建立具备高电荷探测灵敏度、能够对探测空间中异常目标进行实时检测、特征提取与多目标识别能力的传感系统。

在探测系统方面，本研究采用基于MEMS技术的电荷传感阵列系统进行目标探测。该系统依据MEMS结构中的电容效应，将目标在时间与空间上的电荷变化转化为相应电压差信号，并输出为16个耦合通道的时间电压原始数据。该传感阵列基于小信号监测机制，采用直接差分测量策略，而非数字滤波算法，而是通过实验数据与背景数据的对消处理实现信号提取，进而为后续目标识别提供数据基础。

在目标特性建模方面，本研究从预期探测目标在真实水下环境中的电荷分布与电流密度生成机制出发，通过建立控制方程、确定目标物性参数，实现对目标电荷行为的等比模型与缩比模型的物理预测。结合有限元仿真工具（如COMSOL）进行真实场景的数值模拟，并通过实际实验验证仿真结果，系统分析探测数据与理论预期之间的差异。同时通过分析仿真工具中等比模型与缩比模型的理论预测关系，根据真实的缩比物理实验，预测真实远距离情景下的目标特性。随后借助回归分析与机器学习方法进行数据拟合与模型优化，最终构建适用于极化带效应水下目标探测的有效体系。需指出的是，实际目标的电荷分布在时间上呈非平稳特性，在空间中存在显著运动，因此需实现在大范围时空背景中有效分离稀疏目标信号，并完成实时识别与判断。

在数据分析方面，本研究主要依托时域分析、频域分析、时频分析及阵列信号处理等方法。首先获取一端连续时间（约20秒）内各通道的时间–电压序列，进行初步时域分析；随后通过快速傅里叶变换（FFT）获取频域中的幅度谱与相位谱，开展频域分析；进一步联合时频分析方法，以刻画信号的非平稳特性；

数据训练的方法主要是机器学习

\section{研究意义}
% TODO: 草稿中“研究意义”仍为模板占位“本文……”，需补写研究意义。
\section{论文结构}
% TODO: 草稿中“论文结构”仍为模板占位，需按最终章节结构补写。
